\begin{proof}[Proof of \cref{obj:lem:uniform-hybrid-arm-risk}]
Fix parameters and counts satisfying the hypotheses. Write \(\mathbb P\) for their probability measure on the given sample space \(\Omega\).

\begin{enumerate}
\item \textit{Identify the joint count law.}
Define the count space and the selected count vector by
\[
\mathbb N_0=\{0,1,2,\ldots\},\qquad
\mathcal E=(\mathbb N_0^4)^d,\qquad
\mathscr T(\omega)
=
\bigl(J_{aj}(\omega),Z_{aj}(\omega),
K_{aj}(\omega),K_{1-a,j}(\omega)\bigr)_{j=1}^d.
\]
For \(j=1,\ldots,d\), define probability measures
\[
\begin{aligned}
\nu_j={}&
\operatorname{Poi}\!\bigl(t_ps_{aj}(P)\bigr)
\otimes\operatorname{Poi}\!\bigl(uq_{aj}(P)\bigr)\\
&\otimes\operatorname{Poi}\!\bigl(ts_{aj}(P)\bigr)
\otimes\operatorname{Poi}\!\bigl(ts_{1-a,j}(P)\bigr),
\qquad
\nu=\bigotimes_{j=1}^d\nu_j,
\end{aligned}
\]
with coordinates ordered as in \(\mathscr T\).

The selected counts form a subfamily of the stipulated independent family: their indices are distinct, including the two factorial-pool indices because \(a\ne1-a\). Thus, for every
\[
\boldsymbol c=(c_j)_{j=1}^d\in\mathcal E,
\qquad
c_j=(c_j^{\mathrm p},c_j^{\mathrm o},
c_j^{\mathrm f},c_j^{\mathrm{opp}}),
\]
independence gives
\[
\begin{aligned}
\mathbb P\{\mathscr T=\boldsymbol c\}
=\prod_{j=1}^d\bigl[
&\mathbb P\{J_{aj}=c_j^{\mathrm p}\}
\mathbb P\{Z_{aj}=c_j^{\mathrm o}\}\\
&\mathbb P\{K_{aj}=c_j^{\mathrm f}\}
\mathbb P\{K_{1-a,j}=c_j^{\mathrm{opp}}\}
\bigr]
=\nu\{\boldsymbol c\}.
\end{aligned}
\]
Indeed, the event on the left is the intersection of these singleton-coordinate events. Equality on singletons determines measures on the countable space \(\mathcal E\), so
\[
\mathbb P\circ\mathscr T^{-1}=\nu.
\]
% lean: CausalSmith.Stat.AnnotationRarearmFrontier.hybrid_four_count_product_law

\item \textit{Establish the population inputs.}
For the fixed law and arm, set
\begingroup\footnotesize
\[
s_j=s_{aj}(P),\qquad
v_j=s_{1-a,j}(P),\qquad
q_j=q_{aj}(P),\qquad
p_j=p_j(P),\qquad
\bar\mu_j=\mu_{aj}(P),
\qquad
\psi_a=\sum_{j=1}^dp_j\bar\mu_j.
\]
\endgroup
The probability definitions imply
\[
s_j,v_j\ge0,\qquad
p_j=s_j+v_j,\qquad
0\le q_j\le s_j,\qquad
0\le\bar\mu_j\le1,\qquad
\sum_{j=1}^dp_j=1.
\]
If \(s_j=0\), the bound \(0\le q_j\le s_j\) gives \(q_j=0\); if \(s_j>0\), the defining quotient for \(\bar\mu_j\) gives the same identity
\[
s_j\bar\mu_j=q_j
\qquad(j=1,\ldots,d).
\]
By \cref{obj:def:model}, the law satisfies \cref{obj:ass:overlap}. On occupied cells, each arm proportion lies in \([\epsilon,1-\epsilon]\), hence
\[
\epsilon p_j\le s_j.
\]
On a null covariate cell, both arm masses vanish and this inequality also holds. In particular,
\[
s_j=0\quad\Longrightarrow\quad
p_j=v_j=q_j=\bar\mu_j=0.
\]
The corresponding outcome count has mean zero, so both branch contributions vanish almost surely.

The public parameters satisfy
\begingroup\small
\[
S=n\epsilon,\qquad N=n+m,\qquad
0<\epsilon\le1,\qquad
S\le N\epsilon,\qquad
\frac S{32}\le u\epsilon,\qquad
\frac N{64}\le t_p,t,\qquad
\frac13\le\frac t{t_p}.
\]
\endgroup
Consequently \(S,N,u,t_p,t\) are strictly positive.

On \(\mathcal E\), define the coordinate counts by
\[
\bigl(\mathsf J_j(\boldsymbol c),\mathsf Z_j(\boldsymbol c),
\mathsf K_j(\boldsymbol c),\mathsf W_j(\boldsymbol c)\bigr)=c_j.
\]
Define real-valued functions on this entire space by
\[
\begin{aligned}
\mathcal G_j
&=\sum_{h=0}^{L-2}g_h\frac{(\mathsf K_j)_h}{(tB)^h},\\
\mathcal P_j
&=\frac{\mathsf Z_j}{u}
\left(1+\frac{\mathsf W_j\mathcal G_j}{tB}\right),\\
\mathcal H_j
&=\frac{\mathsf Z_j}{u}
\left(1+\frac{\mathsf W_j}{\mathsf K_j+1}\right),\\
\mathcal I_j&=\mathbf1\{\mathsf J_j\le k_0\},\qquad
\mathcal V_j=\mathcal I_j\mathcal P_j+(1-\mathcal I_j)\mathcal H_j.
\end{aligned}
\]
Under \(\nu\), the four coordinate counts in cell \(j\) have independent Poisson laws with means
\[
t_ps_j,\qquad uq_j=us_j\bar\mu_j,\qquad ts_j,\qquad tv_j,
\]
respectively. We write \(\mathbb E_\nu\) and \(\operatorname{Var}_\nu\) for expectation and variance under \(\nu\).
% lean: CausalSmith.Stat.AnnotationRarearmFrontier.uniform_hybrid_arm_risk

\item \textit{Calibrate the degree and localization scale.}
Introduce
\[
\eta_S=\log S,\qquad
F_{\mathrm{grow}}=(2^{24})^L,\qquad
\rho_{\mathrm{pil}}=S^{-20}.
\]
Since \(\eta_S\ge4096\), the definition of \(L\) gives
\[
L\ge4,\qquad
L\le\frac{\eta_S}{1024}<L+1,\qquad
\eta_S<1280L.
\]
Using \(\log2\le1\) and \(\log L\le L\),
\[
\log(F_{\mathrm{grow}}L^3)
=24L\log2+3\log L
\le27L
\le\frac{27}{1024}\eta_S
\le\frac{\eta_S}{2}.
\]
Therefore
\[
F_{\mathrm{grow}}L^3\le\sqrt S,
\qquad
e^{-200000L}\le e^{-20\eta_S}=\rho_{\mathrm{pil}}.
\]
Moreover, \(B>0\), and the minimum in its definition yields
\[
t_pB
=H_0L\frac{t_p}{\min\{t_p,t\}}\ge H_0L,
\qquad
tB
=H_0L\frac{t}{\min\{t_p,t\}}\ge H_0L\ge L.
\]

Define the light and heavy index sets, their light mass, and their light count by
\[
\mathcal J_{\mathrm{light}}=\{j:1\le j\le d,\ s_j\le B\},
\qquad
\mathcal J_{\mathrm{heavy}}=\{j:1\le j\le d,\ s_j>B\},
\]
\[
M_{\mathrm{light}}=\sum_{j\in\mathcal J_{\mathrm{light}}}p_j,
\qquad
c_{\mathrm{light}}=|\mathcal J_{\mathrm{light}}|.
\]
For every light cell, overlap gives
\[
p_j\le\frac{s_j}{\epsilon}\le\frac B\epsilon.
\]
Summing this inequality and using total mass one proves
\[
0\le M_{\mathrm{light}}
\le\min\left\{1,\frac{c_{\mathrm{light}}B}{\epsilon}\right\}
\le\frac{dB}{\epsilon},
\qquad
0\le c_{\mathrm{light}}\le d.
\]
These formulas also apply when either index set is empty, with its sum equal to zero.
% lean: CausalSmith.Stat.AnnotationRarearmFrontier.hybrid_degree_calibration
% lean: CausalSmith.Stat.AnnotationRarearmFrontier.light_cell_mass_bound

\item \textit{Bound the inverse-branch variance.}
Use the reciprocal function and numerical constant supplied by
\cref{obj:lem:poisson-inverse-moments}:
\[
g(k)=\frac1{k+1}\quad(k\in\mathbb N_0),
\qquad
C_0=2^{16},
\qquad
C_{\mathrm{inv}}=2+4C_0=262146.
\]
For each cell define the inverse multiplier
\[
\mathcal W_{\mathrm{inv},j}
=1+\mathsf W_jg(\mathsf K_j).
\]
Independence, the Poisson first and second moments, and
\cref{obj:lem:poisson-inverse-moments} give
\[
\begin{aligned}
\mathbb E_\nu[\mathcal W_{\mathrm{inv},j}^2]
&\le
2+2\bigl((tv_j)^2+tv_j\bigr)
\mathbb E_\nu[g(\mathsf K_j)^2]\\
&\le
2+2C_0\frac{(tv_j)^2+tv_j}{(1+ts_j)^2},
\end{aligned}
\]
where the first inequality uses \((1+x)^2\le2+2x^2\). The variance of the independent product is
\[
\begin{aligned}
\operatorname{Var}_\nu(\mathcal W_{\mathrm{inv},j})
&=tv_j\,\mathbb E_\nu[g(\mathsf K_j)^2]
+(tv_j)^2\operatorname{Var}_\nu(g(\mathsf K_j))\\
&\le C_0\left\{
\frac{tv_j}{(1+ts_j)^2}
+\frac{(tv_j)^2ts_j}{(1+ts_j)^4}
\right\}.
\end{aligned}
\]
Also,
\[
\mathbb E_\nu\!\left[\frac{\mathsf Z_j}{u}\right]=q_j,
\qquad
\mathbb E_\nu\!\left[\left(\frac{\mathsf Z_j}{u}\right)^2\right]
=\frac{q_j}{u}+q_j^2.
\]
Factoring its independent first and second moments therefore gives
\[
\operatorname{Var}_\nu(\mathcal H_j)
=
\frac{q_j}{u}\mathbb E_\nu[\mathcal W_{\mathrm{inv},j}^2]
+q_j^2\operatorname{Var}_\nu(\mathcal W_{\mathrm{inv},j}).
\]

For \(s_j>0\), overlap implies
\[
\frac{v_j^2}{s_j}\le\frac{p_j}{\epsilon},
\qquad
q_j,v_j\le\frac{p_j}{\epsilon}.
\]
Indeed, \(\epsilon v_j\le s_j\), and multiplying by \(v_j\le p_j\) proves the first bound. The elementary inequalities
\[
0\le\frac{ts_j}{1+ts_j}\le1,
\qquad
ts_j\le(1+ts_j)^2
\]
give the four rational bounds
\[
\begin{aligned}
\frac{s_j(tv_j)^2}{(1+ts_j)^2}
&\le\frac{v_j^2}{s_j},&
\frac{s_jtv_j}{(1+ts_j)^2}
&\le v_j,\\
\frac{s_j^2tv_j}{(1+ts_j)^2}
&\le\frac{v_j}{t},&
\frac{s_j^2(tv_j)^2ts_j}{(1+ts_j)^4}
&\le\frac{v_j^2}{ts_j}.
\end{aligned}
\]
For example, the first, third, and fourth left-hand sides equal their respective right-hand sides multiplied by
\((ts_j/(1+ts_j))^2\), \((ts_j/(1+ts_j))^2\), and
\((ts_j/(1+ts_j))^4\). Hence \(q_j\le s_j\) gives
\[
q_j\frac{(tv_j)^2+tv_j}{(1+ts_j)^2}
\le\frac{v_j^2}{s_j}+v_j
\le\frac{2p_j}{\epsilon},
\]
\[
q_j^2\left\{
\frac{tv_j}{(1+ts_j)^2}
+\frac{(tv_j)^2ts_j}{(1+ts_j)^4}
\right\}
\le\frac{v_j}{t}+\frac{v_j^2}{ts_j}
\le\frac{2p_j}{t\epsilon}.
\]
Substitution into the variance identity proves
\[
\operatorname{Var}_\nu(\mathcal H_j)
\le C_{\mathrm{inv}}p_j
\left(\frac1{u\epsilon}+\frac1{t\epsilon}\right).
\]
For \(s_j=0\), both sides vanish by the null-cell identities already established. Summing over all cells gives
\[
\sum_{j=1}^d\operatorname{Var}_\nu(\mathcal H_j)
\le C_{\mathrm{inv}}
\left(\frac1{u\epsilon}+\frac1{t\epsilon}\right).
\]
% lean: CausalSmith.Stat.AnnotationRarearmFrontier.hybrid_canonical_inverse_variance
% lean: CausalSmith.Stat.AnnotationRarearmFrontier.hybrid_canonical_inverse_variance_sum

\item \textit{Bound the light polynomial variances.}
Define, for every cell,
\[
\zeta_j=\frac{s_j}{B},
\qquad
\mathcal W_{\mathrm{pol},j}
=1+\frac{\mathsf W_j\mathcal G_j}{tB}.
\]
The hypotheses of \cref{obj:lem:chebyshev-factorial-certificate} hold because \(L\ge4\), \(t,B>0\), \(\zeta_j\ge0\), and \(tB\ge L\). For a light cell it supplies
\[
\mathbb E_\nu[\mathcal G_j^2]
\le392^L(1+\zeta_j)^{2L}
\le1568^L
\le F_{\mathrm{grow}}.
\]
Independence of the two factorial counts and
\(\mathbb E_\nu[\mathsf W_j^2]=(tv_j)^2+tv_j\) then imply
\[
\begin{aligned}
\mathbb E_\nu[\mathcal W_{\mathrm{pol},j}^2]
&\le
2+2\mathbb E_\nu[\mathcal G_j^2]
\left(\frac{v_j^2}{B^2}+\frac{v_j}{tB^2}\right)\\
&\le
2F_{\mathrm{grow}}
\left(1+\frac{v_j^2}{B^2}+\frac{v_j}{tB^2}\right).
\end{aligned}
\]
The outcome count is independent of this multiplier, so
\[
\mathbb E_\nu[\mathcal P_j^2]
=
\left(\frac{q_j}{u}+q_j^2\right)
\mathbb E_\nu[\mathcal W_{\mathrm{pol},j}^2].
\]

For \(s_j\le B\), the following bounds retain each mixed monomial:
\[
\begin{aligned}
\frac{s_jv_j^2}{B^2}&\le\frac{p_j}{\epsilon},&
\frac{s_jv_j}{tB^2}&\le v_j,&
\frac{s_j^2v_j}{tB^2}&\le\frac{p_j}{t},\\
\frac{s_j^2v_j^2}{B^2}&\le p_j^2,&
s_j^2&\le p_j^2,&
p_j&\le\frac B\epsilon.
\end{aligned}
\]
When \(s_j>0\), the first follows from
\(s_jv_j^2/B^2\le v_j^2/s_j\le p_j/\epsilon\).
The second follows from \(s_j\le B\le tB^2\), since \(tB\ge1\).
The next two use \(s_j^2\le B^2\) and \(v_j\le p_j\).
When \(s_j=0\), all these inequalities follow from \(v_j=p_j=0\).

Expanding the product and using \(q_j\le s_j\) gives
\[
\begin{aligned}
&\left(\frac{q_j}{u}+q_j^2\right)
\left(1+\frac{v_j^2}{B^2}+\frac{v_j}{tB^2}\right)\\
&\quad\le
\frac{s_j}{u}
+\frac{s_jv_j^2}{uB^2}
+\frac{s_jv_j}{utB^2}
+s_j^2+\frac{s_j^2v_j^2}{B^2}
+\frac{s_j^2v_j}{tB^2}\\
&\quad\le
\frac{3p_j}{u\epsilon}+\frac{p_j}{t}+2p_j^2
\le
\frac{3p_j}{u\epsilon}+\frac{p_j}{t}
+\frac{2B^2}{\epsilon^2}.
\end{aligned}
\]
Here \(s_j,v_j\le p_j/\epsilon\) uses \(\epsilon\le1\).
Thus
\[
\operatorname{Var}_\nu(\mathcal P_j)
\le\mathbb E_\nu[\mathcal P_j^2]
\le
6F_{\mathrm{grow}}
\left(\frac{p_j}{u\epsilon}+\frac{p_j}{t}
+\frac{B^2}{\epsilon^2}\right).
\]
Summation yields
\[
\sum_{j\in\mathcal J_{\mathrm{light}}}
\operatorname{Var}_\nu(\mathcal P_j)
\le
6F_{\mathrm{grow}}
\left(
\frac{M_{\mathrm{light}}}{u\epsilon}
+\frac{M_{\mathrm{light}}}{t}
+\frac{c_{\mathrm{light}}B^2}{\epsilon^2}
\right).
\]
% lean: CausalSmith.Stat.AnnotationRarearmFrontier.light_polynomial_variance_sum

\item \textit{Assemble the light mixture variances.}
Define
\[
\pi_j=\nu\{\mathsf J_j\le k_0\}\in[0,1],
\qquad
\delta_j=\mathbb E_\nu[\mathcal P_j]
-\mathbb E_\nu[\mathcal H_j].
\]
The factorial unbiasedness in
\cref{obj:lem:chebyshev-factorial-certificate} gives
\[
\mathbb E_\nu[\mathcal G_j]=G_L(\zeta_j).
\]
For \(s_j>0\), independence and the reciprocal expectation from
\cref{obj:lem:poisson-inverse-moments} therefore give
\[
\begin{aligned}
\mathbb E_\nu[\mathcal P_j]
&=s_j\bar\mu_j
\left(1+\frac{v_j}{B}G_L(\zeta_j)\right)
=p_j\bar\mu_j-v_j\bar\mu_jE_L(\zeta_j),\\
\mathbb E_\nu[\mathcal H_j]
&=s_j\bar\mu_j
\left(1+tv_j\frac{1-e^{-ts_j}}{ts_j}\right)
=p_j\bar\mu_j-v_j\bar\mu_je^{-ts_j}.
\end{aligned}
\]
The polynomial identity used here is
\[
\zeta_jG_L(\zeta_j)=1-E_L(\zeta_j),
\]
as specified in \cref{obj:def:estimator-handle}. The two branch-mean formulas also hold when \(s_j=0\), because \(p_j=v_j=0\).

The pilot is independent of both branches. Consequently,
\[
\begin{aligned}
\mathbb E_\nu[\mathcal V_j]
&=\pi_j\mathbb E_\nu[\mathcal P_j]
+(1-\pi_j)\mathbb E_\nu[\mathcal H_j],\\
\mathbb E_\nu[\mathcal V_j^2]
&=\pi_j\mathbb E_\nu[\mathcal P_j^2]
+(1-\pi_j)\mathbb E_\nu[\mathcal H_j^2].
\end{aligned}
\]
Subtracting the square of the first identity from the second gives
\[
\operatorname{Var}_\nu(\mathcal V_j)
=
\pi_j\operatorname{Var}_\nu(\mathcal P_j)
+(1-\pi_j)\operatorname{Var}_\nu(\mathcal H_j)
+\pi_j(1-\pi_j)\delta_j^2.
\]

For \(0<s_j\le B\), the residual bound in
\cref{obj:lem:chebyshev-factorial-certificate} yields
\[
0\le v_j\bar\mu_jE_L(\zeta_j)
\le\frac{v_jB}{L^2s_j}
\le\frac B{\epsilon L^2}.
\]
The last inequality follows from \(\epsilon v_j\le s_j\). Using the branch means and \((x-y)^2\le2x^2+2y^2\), we obtain
\[
\pi_j(1-\pi_j)\delta_j^2
\le
2\left(\frac B{\epsilon L^2}\right)^2
+2\bigl(p_je^{-ts_j}\bigr)^2.
\]
For \(s_j=0\), the branch means vanish and the same bound applies.

Overlap and the elementary bound \(xe^{-x}\le e^{-1}\) for \(x\ge0\) give
\[
\bigl(p_je^{-ts_j}\bigr)^2
\le p_j^2e^{-2t\epsilon p_j}
\le\frac{p_j}{2e\,t\epsilon}.
\]
Thus
\[
\sum_{j\in\mathcal J_{\mathrm{light}}}
\pi_j(1-\pi_j)\delta_j^2
\le
2c_{\mathrm{light}}
\left(\frac B{\epsilon L^2}\right)^2
+\frac1{e\,t\epsilon}.
\]
Combining this with the polynomial variance bound and
\(0\le\pi_j\le1\) proves
\[
\begin{aligned}
\sum_{j\in\mathcal J_{\mathrm{light}}}
\operatorname{Var}_\nu(\mathcal V_j)
\le{}&
\sum_{j\in\mathcal J_{\mathrm{light}}}
\operatorname{Var}_\nu(\mathcal H_j)\\
&+6F_{\mathrm{grow}}
\left(
\frac{M_{\mathrm{light}}}{u\epsilon}
+\frac{M_{\mathrm{light}}}{t}
+\frac{c_{\mathrm{light}}B^2}{\epsilon^2}
\right)\\
&+2c_{\mathrm{light}}
\left(\frac B{\epsilon L^2}\right)^2
+\frac1{e\,t\epsilon}.
\end{aligned}
\]
% lean: CausalSmith.Stat.AnnotationRarearmFrontier.hybrid_branch_means
% lean: CausalSmith.Stat.AnnotationRarearmFrontier.light_hybrid_variance_sum

\item \textit{Control the heavy mixture variances.}
For a heavy cell define its pilot mean by
\[
\lambda_{\mathrm p,j}=t_ps_j>0.
\]
The cutoff satisfies \(k_0<\lambda_{\mathrm p,j}/4\). The Poisson lower-tail Bernstein inequality states
\[
\nu\left\{
\lambda_{\mathrm p,j}-\mathsf J_j
>\sqrt{2\lambda_{\mathrm p,j}\xi}
\right\}\le e^{-\xi}
\qquad(\xi\ge0).
\]
Since
\[
\sqrt{2\lambda_{\mathrm p,j}(\lambda_{\mathrm p,j}/4)}
=\frac{\lambda_{\mathrm p,j}}{\sqrt2}
<\frac{3\lambda_{\mathrm p,j}}4,
\]
the event \(\{\mathsf J_j\le k_0\}\) is contained in the displayed deviation event with \(\xi=\lambda_{\mathrm p,j}/4\). Hence
\[
\pi_j\le e^{-t_ps_j/4}.
\]

For every \(\zeta_j\ge0\),
\cref{obj:lem:chebyshev-factorial-certificate} supplies
\[
\mathbb E_\nu[\mathcal G_j^2]
\le392^L(1+\zeta_j)^{2L}
\le F_{\mathrm{grow}}(1+\zeta_j)^{2L}.
\]
The multiplier calculation used for light cells now gives
\[
\mathbb E_\nu[\mathcal W_{\mathrm{pol},j}^2]
\le
2F_{\mathrm{grow}}(1+\zeta_j)^{2L}
\left(1+\frac{v_j^2}{B^2}+\frac{v_j}{tB^2}\right).
\]
For \(s_j>B\), one has \(ts_j\ge tB\ge1\). Applying the mixed-monomial calculation with localization scale \(s_j\), while retaining the squared cell mass, gives
\[
\left(\frac{q_j}{u}+q_j^2\right)
\left(1+\frac{v_j^2}{s_j^2}+\frac{v_j}{ts_j^2}\right)
\le
\frac{3p_j}{u\epsilon}+\frac{p_j}{t}+2p_j^2.
\]
Also,
\[
\begin{aligned}
1+\frac{v_j^2}{B^2}+\frac{v_j}{tB^2}
&=
1+\zeta_j^2
\left(\frac{v_j^2}{s_j^2}+\frac{v_j}{ts_j^2}\right)\\
&\le
(1+\zeta_j)^2
\left(1+\frac{v_j^2}{s_j^2}+\frac{v_j}{ts_j^2}\right).
\end{aligned}
\]
Factoring the independent outcome second moment therefore proves
\[
\mathbb E_\nu[\mathcal P_j^2]
\le
6F_{\mathrm{grow}}(1+\zeta_j)^{2L+2}
\left(\frac{p_j}{u\epsilon}+p_j^2+\frac{p_j}{t}\right).
\]

The pilot tail absorbs this growth. Indeed, \(\zeta_j>1\),
\(t_ps_j\ge H_0L\zeta_j\), \(\log(1+\zeta_j)\le\zeta_j\), and
\(2L+2\le(5/2)L\). Thus
\[
\begin{aligned}
&\log\left(
e^{-t_ps_j/4}F_{\mathrm{grow}}(1+\zeta_j)^{2L+2}
\right)\\
&\qquad\le
-\frac{H_0L\zeta_j}{4}
+24L+\frac52L\zeta_j\\
&\qquad\le
-\left(\frac{H_0}{4}-24-\frac52\right)L\zeta_j
\le-200000L.
\end{aligned}
\]
Consequently,
\[
\pi_jF_{\mathrm{grow}}(1+\zeta_j)^{2L+2}
\le e^{-200000L}\le\rho_{\mathrm{pil}},
\qquad
\pi_j\le\rho_{\mathrm{pil}}.
\]
Since \(0\le p_j\le1\) and \(\sum_jp_j=1\),
\[
\sum_{j\in\mathcal J_{\mathrm{heavy}}}
\pi_j\mathbb E_\nu[\mathcal P_j^2]
\le
6\rho_{\mathrm{pil}}
\left(\frac1{u\epsilon}+1+\frac1t\right).
\]

The inverse branch mean lies in \([0,p_j]\), by its explicit formula. Moreover,
\((\mathbb E_\nu[\mathcal P_j])^2\le
\mathbb E_\nu[\mathcal P_j^2]\). Hence
\[
\pi_j(1-\pi_j)\delta_j^2
\le
2\pi_j\mathbb E_\nu[\mathcal P_j^2]+2\pi_jp_j^2,
\]
and summing gives
\[
\sum_{j\in\mathcal J_{\mathrm{heavy}}}
\pi_j(1-\pi_j)\delta_j^2
\le
14\rho_{\mathrm{pil}}
\left(\frac1{u\epsilon}+1+\frac1t\right).
\]
Using the mixture variance identity and
\(\operatorname{Var}_\nu(\mathcal P_j)\le
\mathbb E_\nu[\mathcal P_j^2]\), we conclude
\[
\begin{aligned}
\sum_{j\in\mathcal J_{\mathrm{heavy}}}
\operatorname{Var}_\nu(\mathcal V_j)
\le{}&
\sum_{j\in\mathcal J_{\mathrm{heavy}}}
\operatorname{Var}_\nu(\mathcal H_j)\\
&+20\rho_{\mathrm{pil}}
\left(\frac1{u\epsilon}+1+\frac1t\right).
\end{aligned}
\]
Adding the light and heavy bounds proves
\[
\begin{aligned}
\sum_{j=1}^d\operatorname{Var}_\nu(\mathcal V_j)
\le{}&
\sum_{j=1}^d\operatorname{Var}_\nu(\mathcal H_j)\\
&+6F_{\mathrm{grow}}
\left(
\frac{M_{\mathrm{light}}}{u\epsilon}
+\frac{M_{\mathrm{light}}}{t}
+\frac{c_{\mathrm{light}}B^2}{\epsilon^2}
\right)\\
&+2c_{\mathrm{light}}
\left(\frac B{\epsilon L^2}\right)^2
+\frac1{e\,t\epsilon}\\
&+20\rho_{\mathrm{pil}}
\left(\frac1{u\epsilon}+1+\frac1t\right).
\end{aligned}
\]
% lean: CausalSmith.Stat.AnnotationRarearmFrontier.false_light_absorption
% lean: CausalSmith.Stat.AnnotationRarearmFrontier.heavy_hybrid_variance_sum
% lean: CausalSmith.Stat.AnnotationRarearmFrontier.hybrid_canonical_variance_sum

\item \textit{Bound the selected bias.}
Define the exponential pilot parameter by
\[
\theta_{\mathrm{pil}}=\log\left(1+\frac t{t_p}\right)>0.
\]
The Poisson exponential moment and exponential Markov inequality give
\[
\begin{aligned}
(1-\pi_j)e^{-ts_j}
&\le
e^{-\theta_{\mathrm{pil}}k_0}
\mathbb E_\nu[e^{\theta_{\mathrm{pil}}\mathsf J_j}]
e^{-ts_j}\\
&=
e^{-\theta_{\mathrm{pil}}k_0}
\exp\left\{t_ps_j(e^{\theta_{\mathrm{pil}}}-1)-ts_j\right\}\\
&=e^{-\theta_{\mathrm{pil}}k_0}.
\end{aligned}
\]
This calculation applies to every \(s_j\ge0\). The ratio condition and
\(\log(1+x)\ge2x/(2+x)\) for \(x\ge0\) imply
\[
\theta_{\mathrm{pil}}\ge\log(4/3)\ge\frac27.
\]
The floor inequality \(k_0+1>t_pB/4\), together with \(L\ge4\), gives
\[
k_0\ge\frac{H_0L}{8}.
\]
It follows that
\[
(1-\pi_j)e^{-ts_j}
\le
\exp\left(-\frac{H_0L}{28}\right)
\le e^{-30000L}
\le e^{-20\eta_S}
=\rho_{\mathrm{pil}}.
\]

The branch means yield, for \(s_j>0\),
\[
\mathbb E_\nu[\mathcal V_j]-p_j\bar\mu_j
=
-v_j\bar\mu_j
\left\{\pi_jE_L(\zeta_j)+(1-\pi_j)e^{-ts_j}\right\}.
\]
For \(s_j\le B\), the light residual estimate therefore gives
\[
\left|\mathbb E_\nu[\mathcal V_j]-p_j\bar\mu_j\right|
\le\frac B{\epsilon L^2}+p_j\rho_{\mathrm{pil}}.
\]

For heavy cells, the coefficient bound in
\cref{obj:lem:chebyshev-factorial-certificate} implies
\[
|G_L(\zeta_j)|
\le
\sum_{h=0}^{L-2}|g_h|\zeta_j^h
\le14^L(1+\zeta_j)^L.
\]
The defining polynomial identity gives
\[
\begin{aligned}
|E_L(\zeta_j)|
&=|1-\zeta_jG_L(\zeta_j)|\\
&\le1+\zeta_j14^L(1+\zeta_j)^L\\
&\le2F_{\mathrm{grow}}(1+\zeta_j)^{2L+2}.
\end{aligned}
\]
Here each of the two terms in the preceding sum is bounded by
\(F_{\mathrm{grow}}(1+\zeta_j)^{2L+2}\).
Using \(v_j\bar\mu_j\le p_j\) and the heavy pilot absorption proves
\[
\pi_j|v_j\bar\mu_jE_L(\zeta_j)|
\le2p_j\rho_{\mathrm{pil}}.
\]
Together with the opposite pilot bound, this gives
\[
\left|\mathbb E_\nu[\mathcal V_j]-p_j\bar\mu_j\right|
\le3p_j\rho_{\mathrm{pil}}
\qquad(s_j>B).
\]
For \(s_j=0\), the bias is zero. Applying the triangle inequality cell by cell and summing \(p_j\) proves
\[
\left|
\sum_{j=1}^d
\bigl(\mathbb E_\nu[\mathcal V_j]-p_j\bar\mu_j\bigr)
\right|
\le\frac{dB}{\epsilon L^2}+3\rho_{\mathrm{pil}}.
\]
Since \(S\ge1\), one has \(0\le\rho_{\mathrm{pil}}\le1\), and hence
\[
\begin{aligned}
\left(
\sum_{j=1}^d
\bigl(\mathbb E_\nu[\mathcal V_j]-p_j\bar\mu_j\bigr)
\right)^2
&\le
\frac{2d^2B^2}{\epsilon^2L^4}
+18\rho_{\mathrm{pil}}^2\\
&\le
\frac{2d^2B^2}{\epsilon^2L^4}
+18\rho_{\mathrm{pil}}.
\end{aligned}
\]
% lean: CausalSmith.Stat.AnnotationRarearmFrontier.pilot_false_heavy_absorption
% lean: CausalSmith.Stat.AnnotationRarearmFrontier.hybrid_selected_squared_bias_sum

\item \textit{Form the canonical mean-squared-error envelope.}
Poisson counts have moments of every order, and the inverse denominators are at least one, so all the branch and hybrid statistics have finite second moments. Under the product measure \(\nu\), distinct cells are independent. Therefore the variance and squared-bias decomposition is
\[
\begin{aligned}
\mathbb E_\nu\left[
\left(\sum_{j=1}^d\mathcal V_j-\psi_a\right)^2
\right]
={}&
\sum_{j=1}^d\operatorname{Var}_\nu(\mathcal V_j)\\
&+
\left(
\sum_{j=1}^d
\bigl(\mathbb E_\nu[\mathcal V_j]-p_j\bar\mu_j\bigr)
\right)^2.
\end{aligned}
\]
Substituting the variance and squared-bias bounds gives
\[
\begin{aligned}
\mathbb E_\nu\left[
\left(\sum_{j=1}^d\mathcal V_j-\psi_a\right)^2
\right]
\le{}&
\sum_{j=1}^d\operatorname{Var}_\nu(\mathcal H_j)\\
&+6F_{\mathrm{grow}}
\left(
\frac{M_{\mathrm{light}}}{u\epsilon}
+\frac{M_{\mathrm{light}}}{t}
+\frac{c_{\mathrm{light}}B^2}{\epsilon^2}
\right)\\
&+2c_{\mathrm{light}}
\left(\frac B{\epsilon L^2}\right)^2
+\frac1{e\,t\epsilon}\\
&+20\rho_{\mathrm{pil}}
\left(\frac1{u\epsilon}+1+\frac1t\right)\\
&+\frac{2d^2B^2}{\epsilon^2L^4}
+18\rho_{\mathrm{pil}}.
\end{aligned}
\]
% lean: CausalSmith.Stat.AnnotationRarearmFrontier.hybrid_canonical_mse_decomposition
% lean: CausalSmith.Stat.AnnotationRarearmFrontier.hybrid_canonical_mse_envelope

\item \textit{Normalize the retained terms.}
Define the nonnegative deterministic envelope
\[
\begin{aligned}
\mathscr E_{\mathrm{rate}}
={}&
\frac1S+\frac{d^2B^2}{\epsilon^2L^4}\\
&+F_{\mathrm{grow}}
\left(
\frac{M_{\mathrm{light}}}{S}
+\frac{M_{\mathrm{light}}}{t}
+\frac{dB^2}{\epsilon^2}
\right)
+\frac{dB^2}{\epsilon^2L^4}
+\rho_{\mathrm{pil}}.
\end{aligned}
\]
The intensity bounds, \(S\le N\epsilon\), and \(\epsilon\le1\) imply the first three bounds below; the experiment hypothesis \(S\ge\exp(4096)\ge1\) gives the fourth:
\[
\frac1{u\epsilon}\le\frac{32}{S},
\qquad
\frac1{t\epsilon}\le\frac{64}{S},
\qquad
\frac1t\le\frac{64}{S},
\qquad
\frac1S\le1.
\]
Using \(c_{\mathrm{light}}\le d\), each term of the preceding mean-squared-error envelope satisfies
\[
\begin{aligned}
\sum_j\operatorname{Var}_\nu(\mathcal H_j)
&\le96C_{\mathrm{inv}}\mathscr E_{\mathrm{rate}},\\
6F_{\mathrm{grow}}
\left(
\frac{M_{\mathrm{light}}}{u\epsilon}
+\frac{M_{\mathrm{light}}}{t}
+\frac{c_{\mathrm{light}}B^2}{\epsilon^2}
\right)
&\le192\mathscr E_{\mathrm{rate}},\\
2c_{\mathrm{light}}
\left(\frac B{\epsilon L^2}\right)^2
&\le2\mathscr E_{\mathrm{rate}},\\
\frac1{e\,t\epsilon}
&\le64\mathscr E_{\mathrm{rate}},\\
20\rho_{\mathrm{pil}}
\left(\frac1{u\epsilon}+1+\frac1t\right)
&\le1940\mathscr E_{\mathrm{rate}},\\
\frac{2d^2B^2}{\epsilon^2L^4}
&\le2\mathscr E_{\mathrm{rate}},\\
18\rho_{\mathrm{pil}}
&\le18\mathscr E_{\mathrm{rate}}.
\end{aligned}
\]
For the polynomial term, the first reciprocal bound multiplies
\(M_{\mathrm{light}}/S\) by \(32\), and the other two nonnegative summands are also bounded by \(32\) times their counterparts in \(\mathscr E_{\mathrm{rate}}\).
For the pilot term, the parentheses are at most \(32+1+64=97\).
Adding the seven inequalities, with \(C_{\mathrm{inv}}=262146\), yields
\[
\mathbb E_\nu\left[
\left(\sum_j\mathcal V_j-\psi_a\right)^2
\right]
\le25168234\,\mathscr E_{\mathrm{rate}}.
\]
% lean: CausalSmith.Stat.AnnotationRarearmFrontier.hybrid_canonical_envelope_rate

\item \textit{Absorb the deterministic envelope.}
Define
\[
A_{\mathrm{bw}}=64H_0,
\qquad
\zeta_{\mathrm{rate}}=\frac{d}{N\epsilon L}.
\]
Since \(\min\{t_p,t\}\ge N/64\),
\[
B\le\frac{A_{\mathrm{bw}}L}{N}.
\]
Substituting this bound into the light mass and bandwidth terms gives
\[
M_{\mathrm{light}}
\le A_{\mathrm{bw}}\zeta_{\mathrm{rate}}L^2,
\]
\[
\frac{dB^2}{\epsilon^2}
\le
\frac{A_{\mathrm{bw}}^2\zeta_{\mathrm{rate}}L^3}{N\epsilon}
\le
\frac{A_{\mathrm{bw}}^2\zeta_{\mathrm{rate}}L^3}{S},
\]
\[
\frac{d^2B^2}{\epsilon^2L^4}
\le A_{\mathrm{bw}}^2\zeta_{\mathrm{rate}}^2,
\qquad
\frac{dB^2}{\epsilon^2L^4}
\le\frac{d^2B^2}{\epsilon^2L^4}
\le A_{\mathrm{bw}}^2\zeta_{\mathrm{rate}}^2,
\]
where the last comparison uses \(d\ge1\).

The calibration \(F_{\mathrm{grow}}L^3\le\sqrt S\) and \(L\ge1\) imply
\[
F_{\mathrm{grow}}M_{\mathrm{light}}
\le A_{\mathrm{bw}}\zeta_{\mathrm{rate}}\sqrt S.
\]
Together with \(1/t\le64/S\), the preceding bandwidth bounds give
\[
\begin{aligned}
\frac{F_{\mathrm{grow}}M_{\mathrm{light}}}{S}
&\le\frac{A_{\mathrm{bw}}\zeta_{\mathrm{rate}}}{\sqrt S},\\
\frac{F_{\mathrm{grow}}M_{\mathrm{light}}}{t}
&\le\frac{64A_{\mathrm{bw}}\zeta_{\mathrm{rate}}}{\sqrt S},\\
F_{\mathrm{grow}}\frac{dB^2}{\epsilon^2}
&\le\frac{A_{\mathrm{bw}}^2\zeta_{\mathrm{rate}}}{\sqrt S}.
\end{aligned}
\]
Consequently,
\[
F_{\mathrm{grow}}
\left(
\frac{M_{\mathrm{light}}}{S}
+\frac{M_{\mathrm{light}}}{t}
+\frac{dB^2}{\epsilon^2}
\right)
\le
(65A_{\mathrm{bw}}+A_{\mathrm{bw}}^2)
\frac{\zeta_{\mathrm{rate}}}{\sqrt S}.
\]
The nonnegative square
\((\zeta_{\mathrm{rate}}-1/\sqrt S)^2\) gives
\[
\frac{2\zeta_{\mathrm{rate}}}{\sqrt S}
\le\zeta_{\mathrm{rate}}^2+\frac1S.
\]
Also \(S\ge1\) implies \(\rho_{\mathrm{pil}}=S^{-20}\le S^{-1}\).
Define the numerical constants
\[
C_{\mathrm{abs}}
=
2+2A_{\mathrm{bw}}^2
+\frac{65A_{\mathrm{bw}}+A_{\mathrm{bw}}^2}{2},
\qquad
C=25168234\,C_{\mathrm{abs}}>0.
\]
All the retained terms are therefore bounded as
\[
\mathscr E_{\mathrm{rate}}
\le
C_{\mathrm{abs}}
\left(\frac1S+\zeta_{\mathrm{rate}}^2\right).
\]
We have proved the canonical estimate
\[
\mathbb E_\nu\left[
\left(\sum_{j=1}^d\mathcal V_j-\psi_a\right)^2
\right]
\le
C\left\{
\frac1S+\left(\frac{d}{N\epsilon L}\right)^2
\right\}.
\]
The displayed construction makes \(C\) a fixed numerical constant.
% lean: CausalSmith.Stat.AnnotationRarearmFrontier.hybrid_intermediate_envelope_rate
% lean: CausalSmith.Stat.AnnotationRarearmFrontier.hybrid_canonical_arm_rate

\item \textit{Transfer the canonical estimate to the given counts.}
Define the squared-error function on the count space by
\[
\mathcal R:\mathcal E\longrightarrow[0,\infty),
\qquad
\mathcal R(\boldsymbol c)
=
\left(
\sum_{j=1}^d\mathcal V_j(\boldsymbol c)-\psi_a
\right)^2.
\]
The definitions of the branches give, for every \(\omega\in\Omega\),
\[
\mathcal V_j(\mathscr T(\omega))=V_{aj}(\omega).
\]
The count map is measurable, and every real-valued function on the countable count space is measurable. Using the joint law established at the outset, change of variables therefore gives
\[
\begin{aligned}
\mathbb E\left[
\left(\sum_{j=1}^dV_{aj}-\psi_a\right)^2
\right]
&=
\int_\Omega\mathcal R(\mathscr T(\omega))\,d\mathbb P(\omega)\\
&=
\int_{\mathcal E}\mathcal R(\boldsymbol c)\,d\nu(\boldsymbol c)\\
&\le
C\left\{
\frac1S+\left(\frac{d}{N\epsilon L}\right)^2
\right\}.
\end{aligned}
\]
Finally, \(\psi_a=\sum_jp_j(P)\mu_{aj}(P)\). The same numerical \(C\) thus applies to every admissible choice in the statement.
% lean: CausalSmith.Stat.AnnotationRarearmFrontier.uniform_hybrid_arm_risk
\end{enumerate}
\end{proof}
